\documentclass[11pt,a4paper]{article}
\usepackage[utf8]{inputenc}
\usepackage[T1]{fontenc}
\usepackage{lmodern}
\usepackage{amsmath}
\usepackage[margin=2cm]{geometry}
\usepackage{enumitem}
\pagestyle{plain}
\setlength{\parindent}{0pt}
\setlength{\parskip}{0pt}
\begin{document}
{\Large\bfseries Tarea de hidrácidos y óxidos}
\vspace{12pt}
\begin{enumerate}[label=\textbf{\arabic*.},leftmargin=*,itemsep=11pt,topsep=0pt]
\item \begin{minipage}[t]{\linewidth}
¿Cuál es la fórmula del ácido fluorhídrico?

\vspace{5pt}
\begin{tabular}{@{}p{0.47\linewidth}p{0.47\linewidth}@{}}
a)~$\mathrm{H_{2}F}$ & b)~HF \\
c)~$\mathrm{HF_{2}}$ & d)~FH \\
\end{tabular}
\end{minipage}

\item \begin{minipage}[t]{\linewidth}
¿Cuál es el nombre tradicional del $\mathrm{H_{2}Te}$?

\vspace{5pt}
\begin{tabular}{@{}p{0.47\linewidth}p{0.47\linewidth}@{}}
a)~ácido telurhídrico & b)~telururo ácido \\
c)~ácido telúrico & d)~ácido teluroso \\
\end{tabular}
\end{minipage}

\item \begin{minipage}[t]{\linewidth}
¿Qué nomenclatura no se usa para los hidrácidos?

\vspace{5pt}
\begin{tabular}{@{}p{0.47\linewidth}p{0.47\linewidth}@{}}
a)~sistemática & b)~stock \\
c)~se usan las tres & d)~tradicional \\
\end{tabular}
\end{minipage}

\item \begin{minipage}[t]{\linewidth}
¿Cuál es el nombre sistemático del $\mathrm{Li_{2}O}$?

\vspace{5pt}
\begin{tabular}{@{}p{0.47\linewidth}p{0.47\linewidth}@{}}
a)~monóxido de litio & b)~óxido lítico \\
c)~dióxido de litio & d)~monóxido de dilitio \\
\end{tabular}
\end{minipage}

\item \begin{minipage}[t]{\linewidth}
¿Cuál es la fórmula del heptaóxido de dibromo?

\vspace{5pt}
\begin{tabular}{@{}p{0.47\linewidth}p{0.47\linewidth}@{}}
a)~$\mathrm{Br_{2}O_{5}}$ & b)~$\mathrm{BrO_{7}}$ \\
c)~$\mathrm{Br_{7}O_{2}}$ & d)~$\mathrm{Br_{2}O_{7}}$ \\
\end{tabular}
\end{minipage}

\item \begin{minipage}[t]{\linewidth}
¿Cuál es la fórmula del anhídrido carbonoso?

\vspace{5pt}
\begin{tabular}{@{}p{0.47\linewidth}p{0.47\linewidth}@{}}
a)~$\mathrm{C_{2}O_{2}}$ & b)~$\mathrm{CO_{2}}$ \\
c)~CO & d)~$\mathrm{C_{2}O}$ \\
\end{tabular}
\end{minipage}

\item \begin{minipage}[t]{\linewidth}
¿Cuál fórmula está escrita correctamente?

\vspace{5pt}
\begin{tabular}{@{}p{0.47\linewidth}p{0.47\linewidth}@{}}
a)~$\mathrm{Ca_{2}O_{2}}$ & b)~CaO \\
c)~$\mathrm{Ba_{2}O_{2}}$ & d)~$\mathrm{Mg_{2}O_{2}}$ \\
\end{tabular}
\end{minipage}

\item \begin{minipage}[t]{\linewidth}
¿Cuál es el nombre stock del HgO?

\vspace{5pt}
\begin{tabular}{@{}p{0.47\linewidth}p{0.47\linewidth}@{}}
a)~óxido de mercurio (II) & b)~monóxido de mercurio (II) \\
c)~óxido de mercurio (I) & d)~óxido mercúrico \\
\end{tabular}
\end{minipage}

\item \begin{minipage}[t]{\linewidth}
¿Cuál es el nombre sistemático del $\mathrm{P_{2}O_{3}}$?

\vspace{5pt}
\begin{tabular}{@{}p{0.47\linewidth}p{0.47\linewidth}@{}}
a)~trióxido de difósforo & b)~anhídrido fosfórico \\
c)~óxido de fósforo (III) & d)~dióxido de trifósforo \\
\end{tabular}
\end{minipage}

\item \begin{minipage}[t]{\linewidth}
¿Qué trío de nombres corresponde al $\mathrm{Br_{2}O_{5}}$?

\vspace{5pt}
\begin{tabular}{@{}p{0.47\linewidth}p{0.47\linewidth}@{}}
a)~anhídrido perbrómico, óxido de bromo (V), dióxido de pentabromo & b)~anhídrido bromoso, óxido de bromo (V), pentaóxido de dibromo \\
c)~anhídrido brómico, óxido de bromo (V), pentaóxido de dibromo & d)~anhídrido brómico, óxido de bromo (III), pentaóxido de dibromo \\
\end{tabular}
\end{minipage}

\end{enumerate}
\end{document}
